\section{Introduction}

Security Operations Centres (SOCs) are undergoing a fundamental transformation.
Autonomous AI agents (powered by large language models (LLMs) and equipped
with tool-use capabilities) now perform threat detection, alert triage,
and incident response tasks that were previously the exclusive domain of
human analysts~\cite{apruzzese2023}. Frameworks such as
LangGraph~\cite{langgraph2024} and AutoGen~\cite{wu2023autogen} have made it
straightforward to deploy multi-step reasoning agents that query SIEM systems,
correlate indicators of compromise (IOCs), and recommend containment actions
with minimal human oversight. Commercial platforms including CrowdStrike
Charlotte AI and SentinelOne Purple AI have begun embedding agent-based
automation into production SOC workflows, and the trend is accelerating.

Yet a critical question remains unanswered: \textit{when do these agents fail,
and how dangerous are those failures?} The security community has focused
heavily on measuring what AI agents can do (detection accuracy, response
speed, ATT\&CK coverage) while paying far less attention to the conditions
under which agents produce incorrect, incomplete, or actively misleading
outputs. This gap carries serious operational consequences:

\begin{itemize}[noitemsep, leftmargin=*]
  \item A \textbf{false negative} allows an intrusion to persist undetected
    for hours or days, expanding the attacker's foothold while the SOC
    believes the environment is clean.
  \item A \textbf{hallucinated threat attribution} propagates incorrect IOCs
    into shared intelligence feeds, poisoning downstream defences across
    multiple organisations that consume the same threat data.
  \item A \textbf{delayed response}, caused by an agent paralysed by missing
    tool access or overwhelmed by input complexity, allows an incident to
    escalate beyond regulatory reporting thresholds.
\end{itemize}

These are not hypothetical scenarios: they are predictable consequences of
deploying agents whose capabilities do not match the demands of their assigned
tasks. We call this the \textit{Skill--Knowledge--Capability (SKC) mismatch}
problem: the systematic gap between what a security task requires and what
an agent can actually deliver across multiple dimensions including procedural
skills, declarative threat knowledge, tool access, reasoning quality, and
knowledge currency.

Existing work addresses fragments of this problem in isolation. Classical
dependability theory~\cite{avizienis2004} predates LLM agents and has no
notion of epistemic uncertainty or knowledge currency.
Hallucination research~\cite{llmagent_halluc_survey2025} treats reasoning
failures as model-internal properties without modelling tool availability,
task complexity, or knowledge freshness. Calibration
work~\cite{guo2017calibration} shows neural networks are miscalibrated but
does not connect this to agent reliability. Agent
benchmarks~\cite{liu2023agentbench, deng2023pentest} measure capability
without modelling failure conditions. The closest
recent efforts —
MAST~\cite{cemri2025mast}, a domain-agnostic 14-mode failure
taxonomy, and AgentDojo~\cite{debenedetti2024agentdojo}, an
adversarial prompt-injection environment — each cover one slice of
the surface. To the best of our knowledge, no existing work provides a
unified reliability model, evaluation testbed, and labelled failure
dataset for the full non-adversarial failure surface of AI security
agents.

\textbf{This paper presents \ARES{}} (Agentic Reliability Evaluation for
Security systems), a framework designed to fill this gap. \ARES{} provides:
\begin{enumerate}[noitemsep, leftmargin=*]
  \item A \textbf{three-axis failure taxonomy} classifying AI security agent
    failures by origin (8 classes), manifestation (8 classes), and operational
    impact (7 classes), grounded in CSIRT practice.
  \item A \textbf{quantitative reliability model} based on SKC mismatch,
    including an Agent Reliability Score (ARS) with task-class-specific
    weights and an Epistemic Gap metric that formalises the hallucination
    precondition.
  \item \textbf{\ARESBench{}}, a fully open-source, Docker-deployable
    evaluation testbed supporting four agent architectures across five
    threat scenario classes.
  \item \textbf{\SAFK{}}, a labelled dataset of 1,000 agent execution
    traces under controlled failure conditions, released under CC~BY~4.0.
\end{enumerate}

Together they let SOC architects screen an agent for a given task before
it ever sees production traffic.

Parallel 2025--2026 efforts confirm the urgency: hallucination
benchmarks~\cite{bang2025hallulens} treat confident factual errors as
a cross-domain failure class, and the OWASP Top~10 for Agentic
Applications~\cite{owasp2026agentic} elevates capability-mismatch and
tool-misuse failures to first-class operational risks. None of these
efforts connect agent reliability to SOC-specific consequences: a
misclassified ransomware beacon or a fabricated CVE attribution.
\ARES{} fills that gap.
